% Build: latexmk -pdf to_do.tex

\documentclass[12pt,a4paper]{article}

\usepackage[T1]{fontenc}
\usepackage[utf8]{inputenc}
\usepackage[english,french]{babel}
\usepackage{amsmath,amssymb,amsthm,mathtools}
\usepackage{microtype}
\usepackage{geometry}
\usepackage{hyperref}
\usepackage{mathptmx}

\geometry{margin=2.5cm}

% Custom commands and theorem environments.

\newcommand{\N}{\mathbb{N}}
\newcommand{\Z}{\mathbb{Z}}
\newcommand{\R}{\mathbb{R}}

\theoremstyle{plain}
\newtheorem{theorem}{Theorem}[section]
\newtheorem{lemma}[theorem]{Lemma}
\newtheorem{proposition}[theorem]{Proposition}
\newtheorem{corollary}[theorem]{Corollary}

\theoremstyle{definition}
\newtheorem{definition}[theorem]{Definition}
\newtheorem{example}[theorem]{Example}

\theoremstyle{remark}
\newtheorem{remark}[theorem]{Remark}


\title{Stage MIDS --- À faire}
\author{}
\date{\today}

\begin{document}
\maketitle

\section{Notations}

\begin{itemize}
    \item $X^w_t$ = nombre total d'occurrences du mot $w$ à l'instant $t$
    \item $N_t$ = nombre total de mots dans le corpus à l'instant $t$
    \item $f_t = \frac{X^w_t}{N_t}$ = fréquence du mot $w$ à l'instant $t$
\end{itemize}

Le problème avec $f_t$ est que ce n'est pas un entier ; je préférerais conserver la structure entière de $X_t$, mais il faut tout de même une forme de normalisation par rapport à $N_t$. Pour cela, je propose de normaliser la moyenne de n'importe quel modèle probabiliste par rapport à $N_t$.

\section{Détection de pics (spikes)}

\subsection{Étape 1 : ajuster un modèle binomial négatif}

La loi binomiale négative de paramètres $p \in (0,1)$, $r > 0$ est définie par
$$\mathbb{P}(X = k) = \frac{\Gamma(k + r)}{\Gamma(k + 1) \Gamma(r)} p^k (1-p)^r. $$
Sa moyenne et sa variance valent
$$
\mathbb{E}[X] = \frac{r(1-p)}{p} \quad \text{et} \quad \text{Var}[X] = \frac{r(1-p)}{p^2}.
$$
Une paramétrisation courante consiste à poser $m = r(1-p)/p$, de sorte que $\mathbb{E}[X] = m$ et $\mathrm{Var}(X) = m + m^2/r$. C'est la paramétrisation que nous utiliserons, c'est-à-dire
$$\mathbb{P}(X = k) = \frac{\Gamma(k + r)}{\Gamma(k + 1) \Gamma(r)} \left(\frac{r}{m+r}\right)^k \left(\frac{m}{m+r}\right)^r.$$

\begin{itemize}
    \item D'après GPT, ajuster une binomiale négative avec cette paramétrisation peut se faire avec \texttt{statsmodels.discrete.discrete\_model.NegativeBinomial}, avec $\phi = r^{-1}$.
    \item Les queues sont exponentielles, i.e. $p_n \approx e^{-c n}$ pour un certain $c$. Les queues deviennent plus légères quand $r$ augmente.
    \item Quand $r \to \infty$, la loi converge vers une loi de Poisson de paramètre $m$ ; ce modèle est donc très polyvalent, absorbant les lois de Poisson, géométrique et de Pascal.
    \item Comme $X^w_t$ peut fortement dépendre du nombre de mots $N_t$, il est naturel de normaliser le modèle linéairement avec $N_t$, i.e. de paramétrer $X_t \sim \mathrm{NB}(\mu N_t, r)$. Il est apparemment possible d'ajuster ce modèle sur les observations $(X_t, N_t)$ à l'aide du paramètre \texttt{exposure} de \texttt{statsmodels.discrete.discrete\_model.NegativeBinomial}. Dans ce cas, $\mu$ sera la fréquence moyenne du mot $w$.
    \item Une fois le modèle ajusté et les paramètres optimaux $\mu_*, r_*$ trouvés, il est important de tracer la densité de $\mathrm{NB}(\mu_* \times 100000, r_*)$ par-dessus l'histogramme des données (le graphe de droite sur les graphes d'Elias) pour voir si l'ajustement est bon ou non.
\end{itemize}

\bigskip

Diagnostics pour la qualité de l'ajustement ? Je ne sais pas. Vérifier peut-être la moyenne, la variance, l'asymétrie (skewness) et l'aplatissement (kurtosis) ?

\subsection{Étape 2 : détecter les pics}

Une fois un modèle ajusté, disons $\mu_*, r_*$, on peut détecter les valeurs aberrantes (outliers) dans les données. Notons $P_t$ la densité de probabilité de $\mathrm{NB}(\mu_* N_t, r_*)$. On pose
$$p_t = P_t([X_t, +\infty]), $$
qui est la $p$-valeur de $X_t$ sous le modèle ajusté. On déclare $X_t$ comme aberrant si $p_t$ est trop petit, disons $p_t < 0{,}0001$.

Remarques.
\begin{itemize}
\item $\log_{10}(p_t)$ est appelé la « surprise » de l'observation $X_t$. Plus elle est élevée, plus l'observation est surprenante.
\item on peut procéder en plusieurs tours : ajuster un premier modèle NB2, puis retirer tous les « outliers évidents » (ceux ayant une surprise très élevée), puis réajuster un modèle sur les données restantes, retirer à nouveau les outliers, etc.
\item Diag : tracer l'histogramme des $p_t$, il devrait être à peu près uniforme, sauf pour quelques outliers trop proches de $0$ ?
\end{itemize}

\subsection{Étape 3 : clusters d'activité anormale}

\subsection{Étape 4 : séries temporelles ?}

Si ce qui précède ne fonctionne pas très bien, on pourrait ajouter un modèle autorégressif sur les paramètres. @Elias : il serait bien de tracer la fonction d'autocorrélation de $X_t$, ou plutôt de $f_t$ ; je m'attends fortement à observer des autocorrélations intéressantes dans $f_t$.

On ajusterait typiquement quelque chose comme $m_t = N_t \times e^{R_t}$ avec $R_t$ un processus AR(1) ?

\newpage
\section{Nouveautés à intégrer (par ordre de difficulté)}

Cette section rassemble les pistes issues des échanges les plus récents,
classées de la plus simple à la plus difficile à mettre en œuvre.

\subsection*{Facile}

\begin{itemize}
    \item \textbf{Commencer par une loi de Poisson.} Avant la binomiale négative,
    ajuster mot par mot un modèle de Poisson, en veillant à faire scaler la moyenne
    avec la taille du corpus : fitter $\mathrm{Poisson}(\lambda N_t)$ et \emph{non}
    $\mathrm{Poisson}(\lambda)$ (c'est la remarque de Benoît sur la normalisation par $N_t$).

    \item \textbf{Batterie d'indicateurs et de graphes, pour 5 ou 6 mots.} Pour chaque mot,
    produire :
    \begin{itemize}
        \item moyenne / écart-type / asymétrie (skewness) / aplatissement (kurtosis)
        de la \emph{série observée} ;
        \item les mêmes moments pour la \emph{loi ajustée} ;
        \item le graphe histogramme des données \emph{vs} densité de la loi ajustée (superposés) ;
        \item l'histogramme des $p$-valeurs de chaque observation.
    \end{itemize}
    Ces sorties suffisent déjà à bien réfléchir à la suite.

    \item \textbf{Rappel sur la $p$-valeur.} Une fois le modèle ajusté, si $F$ est sa
    fonction de répartition, la $p$-valeur d'une observation $X$ est simplement
    $1 - F(X) = \mathbb{P}(\text{être} > X)$. Une $p$-valeur trop petite $\Rightarrow$ outlier.
\end{itemize}

\subsection*{Moyen}

\begin{itemize}
    \item \textbf{Regarder le différentiel de fréquences $X_t - X_{t-1}$}, et pas seulement
    $X_t$. Ce différentiel peut être négatif : on obtient donc des sauts vers le haut
    (quantile supérieur) \emph{et} vers le bas (quantile inférieur). Un saut anormal
    \emph{vers le bas} peut s'interpréter comme une volonté de censure ou d'invisibilisation.
    Il faudra comparer ce signal à la fréquence elle-même.

    \item \textbf{Distinguer soigneusement le différentiel de la fréquence.}
    Difficulté observée : les pics négatifs semblent souvent n'être que le
    contrecoup d'un pic positif précédent (retour à la normale après une hausse),
    plutôt qu'une vraie censure. Deux pistes complémentaires :
    \begin{itemize}
        \item passer à un modèle autorégressif d'ordre supérieur, AR(2) ou AR(3),
        pour absorber cet effet de rebond ;
        \item garder à l'esprit qu'il est aussi possible qu'il n'y ait tout simplement
        pas de phénomène de censure sur les mots regardés.
    \end{itemize}

    \item \textbf{Rendre le choix du seuil « principled » plutôt qu'arbitraire.}
    Le seuil actuel (99{,}9\,\% ou 99\,\%) est arbitraire ; selon les mots, 97\,\%
    pourrait suffire. Objectif : déduire le seuil des données. Pistes :
    \begin{itemize}
        \item identifier le point où l'ajustement d'une loi devient « assez mauvais »
        (le bulk lisse de la distribution \emph{vs} les points qui n'en font pas partie) ;
        \item transposer un critère d'outlier type distance de Cook, en l'adaptant
        aux données discrètes (comptages) ;
        \item construire un test statistique : estimer la loi \emph{en éliminant} les
        outliers, puis tester si les occurrences suspectes sont « naturelles » sous cette loi.
    \end{itemize}

    \item \textbf{Cadrer outliers vs breakpoints.} Distinction à garder claire :
    pour l'instant on cherche des \emph{outliers} (périodes courtes d'activité anormale :
    un jour, une semaine, un mois — typiquement le Covid), car c'est le plus simple et
    probablement déjà très significatif. La question des \emph{breakpoints} (changement
    total et définitif de toute la distribution, typiquement un rachat comme
    Bolloré $\to$ JDD) viendra dans un second temps.
\end{itemize}

\subsection*{Difficile}

\begin{itemize}
    \item \textbf{Ajuster une loi commune à plusieurs mots.} Au lieu d'ajuster
    les occurrences d'un seul mot, ajuster simultanément les occurrences de tout
    (ou d'une partie) du vocabulaire, pour disposer d'un critère d'outlier partagé.
    Plus difficile que l'approche mot par mot ; à tenter dans un second temps.

    \item \textbf{Détecter les breakpoints (ruptures durables).} Au-delà des pics
    ponctuels, détecter un déplacement définitif de la distribution après un événement
    (rachat de journal). Suppose des outils de détection de rupture, pas seulement
    de détection d'outliers.

    \item \textbf{Question ouverte : le bruit décroît-il vraiment en $\sqrt{n}$ en
    lexicométrie ?} Intuition de Benoît : probablement \emph{non} (comportement
    non-diffusif). Deux raisons avancées :
    \begin{itemize}
        \item les mots ne sont pas des variables indépendantes : le texte est un tissu ;
        \item la taille d'échantillon effective (\emph{effective sample size}) est très
        inférieure à la taille brute $n$ (nombre de tokens) : on ne peut pas répéter
        « inflation » trois fois par phrase (facteur constant), et le texte se répète
        (effet non linéaire, ESS probablement concave en $n$).
    \end{itemize}
    Conséquence pratique : on observe des valeurs plus extrêmes sur un petit corpus
    (un hebdo) que sur un grand (un quotidien). Le modèle standard « mots i.i.d. »
    donnerait une loi binomiale à variance décroissant en $\sqrt{n}$ ; il n'existe pas
    de littérature sur la construction d'intervalles de confiance en lexicométrie.
    Piste suggérée : du \textbf{bootstrap} pour mesurer empiriquement le rythme de
    convergence.
\end{itemize}

\end{document}
